% ============================================================================
%  qicert — Sample Technical Report (Global Quantum + AI Challenge 2026)
%  Format: IEEE conference style, two columns, 4-8 pages (challenge requirement)
%  Built from submission/10-report-outline.md. Result tables are PRE-REGISTERED
%  (targets + placeholders); measured values land from the experiment ledger
%  (research/experiments.md) during weeks 1-6. No fabricated numbers.
% ============================================================================
\documentclass[conference,10pt]{IEEEtran}

\usepackage{amsmath,amssymb,amsthm}
\usepackage{booktabs}
\usepackage{graphicx}
\usepackage{array}
\usepackage{url}
\urlstyle{same}

\newtheorem{theorem}{Theorem}

\newcommand{\F}{\widetilde{F}}
\newcommand{\norm}[1]{\left\lVert #1 \right\rVert}

\begin{document}

\title{qicert: Certified Quantum-Inspired Compression of Vision--Language--Action
Models with Exact Bounds, Provable Pruning, and Formal Safety Tails}

\author{\IEEEauthorblockN{The qicert Team}
\IEEEauthorblockA{Global Quantum + AI Challenge 2026 --- Autonomous Driving \& Robotics Tracks\\
Code: \url{https://github.com/ZakLr/qicert} \quad Reproduces every table via \texttt{qicert.bench.all}}}

\maketitle

% ---------------------------------------------------------------------------
\begin{abstract}
Vision--Language--Action Models (VLAMs) at 7B--10B parameters are too large for
embedded vehicle and robot controllers and too expensive to certify. We present
\textbf{qicert}, a backbone-agnostic pipeline that compresses VLAMs into
quantized tensor-train (QTT) networks before the heaviest training, compiles
cross-modal interaction blocks into exactly-evaluable commuting-Pauli families,
and certifies what remains with a three-layer certificate stack: exact
algebraic bounds (Lipschitz products of QTT cores; per-family pruning errors),
local formal proofs (SOS boxes; Lyapunov-margin survival curves), and
statistical tail bounds on signal-temporal-logic robustness estimated by
iterative quantum amplitude estimation raced against the strongest honest
classical estimator (RESTART+GEV) and closed by scenario-optimization and
conformal calibration. A syndrome--shadow monitor guards the action stream with
theorem-budgeted false-alarm rates. We report the honest measured status on a
fine-tuned MiniVLA-1B (full frozen 108-episode split, validated harness,
FT-anchor exactness $+0.0000$): calibrated INT8 is accuracy leader at
2$\times$ compression (0.4466 vs 0.4468 dense), while our certified
TT+residual track measurably trails on accuracy at 2.86$\times$ (0.1529) ---
and is the only compressed model in the comparison that ships sound runtime
certificates, enforced by a load/step guard that refuses every action or
artifact it cannot certify. The INT8 baseline cannot compute these
certificates, which are exact functions of the compressed cores.
\end{abstract}

\begin{IEEEkeywords}
tensor networks, quantized tensor train, Lipschitz certificates, Pauli grouping,
Clifford diagonalization, quantum amplitude estimation, signal temporal logic,
classical shadows, VLAM, autonomous driving, robotics, formal safety
\end{IEEEkeywords}

% ---------------------------------------------------------------------------
\section{Introduction}

VLAMs fuse camera, LiDAR, radar, and language inputs into planning and control
outputs. Reference models such as Alpamayo-R1~\cite{alpamayo} and
$\pi_0$~\cite{pi0} demonstrate the category's potential, but their scale makes
them prohibitively expensive to train and nearly impossible to deploy on
embedded hardware within the 100\,ms real-time budget while satisfying
functional-safety certification (ISO 26262; IEC 62061 / ISO 10218).

The challenge asks for a reproducible quantum-enhanced or quantum-inspired
component with a clear advantage over a classical baseline across four
bottlenecks: model footprint, attention complexity, RL alignment cost, and
control-safety guarantees. Quantum-inspired tensor networks have demonstrated
strong LLM compression \cite{compactifai}, and provably efficient quantum
algorithms promise gains in the high-dimensional linear algebra central to VLAM
training \cite{liu}. Our answer is a \emph{system}: one pipeline in which
five quantum-inspired primitives --- tensor networks, Pauli observables,
amplitude encoding, syndrome measurement, and shadow tomography --- are the
type system of every stage (train $\rightarrow$ compress $\rightarrow$ compile
$\rightarrow$ certify $\rightarrow$ monitor). At the same time, the
\emph{classical-energy} tier (a MiniVLA-1B robotics baseline on the LIBERO
spatial slice, with INT8 weight-only quantization + claws) is the
cross-cutting fidelity and energy context for Q3: it tells us, honestly,
whether the classical-quantized relationship is actually the interesting
physics or whether quantum matters. The two tiers share one repository and one
ledger but have separate hardware, separate stacks, and separate kill
criteria; they are not the same experiment presented twice. The contributions are:

\begin{enumerate}
\item \textbf{Exact-certificate compression}: QTT cores double as the objects
of an exact global Lipschitz bound --- exact per-layer factors, residual gap
measured (Sec.~\ref{sec:cert}) --- giving a certified-safe-set column that no
quantization-based baseline can produce at any compression ratio.
\item \textbf{Compilation-as-identity}: cross-modal interaction tensors are
grouped into commuting Pauli families, simultaneously diagonalized by single
Cliffords, and evaluated exactly; dropping a family costs an exact, printed
error $\norm{\Sigma_{a\in F_j} c_a}_1$ (Sec.~\ref{sec:compiler}).
\item \textbf{Quadratic query-complexity savings on the safety tail (realized
oracle)}: the certified failure probability is estimated by iterative quantum
amplitude estimation (IQAE) with Bayesian credible intervals, over a
structured scenario domain with a closed-form certificate-margin oracle
(Sec.~\ref{sec:safety}), raced end-to-end against RESTART+GEV at matched total
cost and closed by Campi--Garatti scenario-optimization and conformal
prediction.
\item \textbf{A runnable, self-reproducing artifact}: \texttt{qicert}
(``certified quantum-inspired compression'') reproduces every table from a clean
environment and emits its own hardware pathway (qubits, depth, shots).
\end{enumerate}

All quantum components are simulated per challenge constraints (Sec.~5.3);
resource assumptions are stated per algorithm. Results below are reported as
pre-registered targets with measured values pending (Sec.~\ref{sec:results}).

% ---------------------------------------------------------------------------
\section{System Architecture}
\label{sec:arch}

The pipeline is identical across both tracks; only the slice specification
differs.

\textbf{Training time (offline, Kaggle 2$\times$T4):}
\textbf{Stage 0} baseline fine-tune (QLoRA/LoRA) of the backbone
(AD: LLaVA-1.5-7B with Alpamayo-adjacent action head, justified per Sec.~5.4;
robotics: MiniVLA-class head). \textbf{Stage 1} compress into $\{$TT, QTT$\}$
via TT-cross (query-based, no global SVD); the safety-budgeted bond allocator
sets per-layer ranks. \textbf{Stage 2} ALS/DMRG-native fine-tuning on the
compressed manifold --- the dense matrix is never formed. \textbf{Stage 3}
compile cross-modal interaction blocks into commuting-Pauli families.
\textbf{Stage 4} certify (three layers, Sec.~\ref{sec:cert}). \textbf{Stage 5}
package into \texttt{qicert}; the bench suite reproduces every table.

\textbf{Runtime (deployed, $\leq$100\,ms profile):}
sensor + instruction encoders $\rightarrow$ compiled interaction kernels
(diagonal, $O(N\cdot k)$-style) $\rightarrow$ QTT-compressed action head with
certified margins $\rightarrow$ syndrome--shadow monitor (PDU hygiene + shadow
statistics) $\rightarrow$ safe output or certified fallback action.

\textbf{Fusion joints} (why this is a system, not a portfolio):
\begin{table}[!t]
\caption{The seven fusion joints. Each component consumes another's output, so
ablations isolate cleanly and the marginal cost of every extra claim is near zero.}
\label{tab:joints}
\centering
\small
\begin{tabular}{@{}cl@{}}
\toprule
\textbf{Joint} & \textbf{What it buys} \\
\midrule
1 & QTT cores = certificate inputs (zero extra compute) \\
2 & Compiled families = monitor's measurement basis (zero extra design) \\
3 & Certificate margins = the STL predicates IQAE estimates (compositional) \\
4 & Sec.~5.2 seeded scenarios = the amplitude-estimation oracle \\
5 & Worst-case long-tail margin = bond-allocation objective (tail designs model) \\
6 & Compiled qubit/depth/shots table = the hardware pathway \\
7 & \texttt{qicert} carries everything \\
\bottomrule
\end{tabular}
\end{table}

\begin{figure}[!t]
\centering
\fbox{\parbox{0.92\linewidth}{\small \centering
Stage 0: baseline fine-tune (QLoRA/LoRA) \\[1mm]
$\downarrow$ Stage 1: $\{$TT,QTT$\}$ TT-cross compression + bond allocator \\[1mm]
$\downarrow$ Stage 2: ALS/DMRG-native fine-tuning (dense matrix never formed) \\[1mm]
$\downarrow$ Stage 3: compile commuting-Pauli families (exact, prunable) \\[1mm]
$\downarrow$ Stage 4: certify --- L1 exact Lipschitz / L2 SOS + Lyapunov / L3 IQAE tails \\[1mm]
$\downarrow$ Stage 5: package \texttt{qicert}; bench suite reproduces every table}}
\caption{Training-time pipeline (Sec.~\ref{sec:arch}); the runtime view adds the
syndrome--shadow monitor (Sec.~\ref{sec:monitor}).}
\label{fig:pipeline}
\end{figure}

\begin{center}
\fbox{\parbox{0.94\linewidth}{\small \textbf{Unifying principle.} A VLAM that will
be certified and deployed on embedded hardware is, at every stage, a computation
on \emph{structured tensors} --- weights, interaction tensors, belief and
robustness distributions. Quantum information is the discipline that made
structured tensors a first-class programming model: tensor networks, Pauli
observables, amplitude encoding, syndrome measurement, shadow tomography. We
build the entire train $\rightarrow$ compress $\rightarrow$ compile
$\rightarrow$ certify $\rightarrow$ monitor pipeline from those five
primitives. Quantum-inspiration is the type system of the whole submission, not
a label on one component.}}
\end{center}

% ---------------------------------------------------------------------------
\section{The Certificate Stack}
\label{sec:cert}

Three layers answer three questions and are never confused.

\subsection{Layer 1: Exact algebraic bounds}

For a TT/QTT layer with cores $G_1\dots G_d$ the reconstructed weight matrix
is $W = \mathrm{contract}(G_1,\dots,G_d)$, and norm submultiplicativity gives
\begin{equation}
\norm{W}_2 \;\le\; \prod_{i=1}^{d} \norm{G_i}_2 ,
\label{eq:lipschitz}
\end{equation}
each core norm computed by power iteration on the small core in
$O(\sum_i r_i^2 n_i)$. The bound is deliberately conservative; the
\emph{exact} per-layer operator norm is available at the same cost class:
power iteration on the TT contraction performs matvecs $x \mapsto Wx$ in
$O(\sum_i r_{i-1} n_i^2 r_i)$ per step (standard TT contraction order;
$O(d n^2 r^2)$ at uniform mode size) without ever forming $W$. Chaining through the
network with 1-Lipschitz activations (ReLU et al.),
\begin{equation}
L(\F) \;=\; \prod_{\text{layers }l} \norm{W^{(l)}}_2 ,
\end{equation}
where each factor is the exact per-layer norm (power iteration), not the
core-product bound. The residual gap between $L(\F)$ and
$\sup_x \norm{\nabla \F(x)}_2$ --- activation interleaving and non-orthogonality
across layers --- is \emph{measured}, not assumed: N3 registers the tightness
ratio $\kappa = L(\F)/\widehat{L}_\nabla$, where $\widehat{L}_\nabla$ estimates
$\sup_{x\in B}\norm{\nabla \F(x)}_2$ by power iteration on the full TT network
at adversarial sample points in $B$, and prints it at every Pareto point.
With the input box $B$ from the slice spec this yields the \emph{exact
safe-set certificate}
\begin{equation}
\mathrm{Safe}(B) \;=\; \big\{ x \in B : \norm{\F(x)-\F(x_{\mathrm{ref}})} \le
L(\F)\,\mathrm{diam}(B) < \mathrm{margin} \big\}.
\label{eq:safeset}
\end{equation}
Non-emptiness is a \emph{measured} column (Table~\ref{tab:main}), never an
assumption: a conservative $L$ shrinks the column but cannot falsify it, and
the column still separates us from INT8, for which no certificate object
exists at any ratio. Its usefulness is gated by R2; where the global bound is
loose, Layer~2a supplies the tight local boxes.

\begin{theorem}[Lipschitz product, layer 1]
For a TT layer with cores $G_1\dots G_d$, Eq.~\eqref{eq:lipschitz} holds with
no sampling and no input-distribution assumption, and the tighter exact factor
$\norm{W}_2$ is computable by power iteration on the TT contraction at cost
$O(k_{\mathrm{it}}\sum_i r_{i-1} n_i^2 r_i)$ with $k_{\mathrm{it}}$ iterations,
the dense matrix never formed. Both properties are preserved for QTT (log-mode)
reshapes because they depend only on the contraction, not on mode sizes.
\end{theorem}

The load-bearing comparison: INT8 quantization error is data-dependent
(clip ranges, per-channel scales) and unfactorizable --- no exact global
constant exists. SVD truncation gives no certificate object to compose. Our
compressed cores \emph{are} the certificate.

\begin{figure}[!t]
\centering
\fbox{\parbox{0.92\linewidth}{\small \centering
[Measured in N3/N5: certified-safe-set \% vs compression ratio] \\[1mm]
qicert-QTT: monotone decay, bond-spectrum predictor band \\
INT8 / SVD: not computable (flat 0\% line, stated honestly)}}
\caption{The asymmetry plot (pre-registered): the certificate column that
quantization baselines structurally cannot produce.}
\label{fig:safeset}
\end{figure}

\subsection{Layer 2: Local formal proofs}

\textbf{(a) SOS-certified boxes.} Composition polynomials of TT-structured
networks admit moment-SOS hierarchies with tractable SDPs~\cite{sos-tt}.
SOS is deliberately \emph{local by construction}: each certificate covers one
sampled long-tail box $B_k$ and a low-degree ($\le 4$) composition polynomial
over a small slice (state dimension $\lesssim 16$), never a whole-network SDP
--- formal verification of a 7B network in one SDP is out of scope for any
method, and the hierarchy's state-lifting structure keeps each box small.
R3 budgets each box at $<$2 h and demotes SOS to a report line if violated.
For each sampled long-tail box $B_k$ we certify
\begin{equation}
\forall x \in B_k:\;\; \mathrm{margin}(x) \ge m_k \ge 0,
\end{equation}
proven by an exported SOS certificate (machine-checkable Gram matrix).
Layer 1 is global-but-loose; Layer 2a is local-but-tight; disagreement between
them is a surfaced bug, not a hidden number.

\textbf{(b) Lyapunov-margin degradation curve + predictor.} The old binary
``certificate survives'' claim is upgraded to a continuous object. At each
compression ratio $r$,
\begin{equation}
\mathrm{survival}(r) \;=\; \frac{\#\{\text{trajectories where } V\text{-decrease holds}\}}
{\#\{\text{trajectories}\}},
\end{equation}
with a regression from (bond spectra, entropy signals, layer sensitivities) to
$\mathrm{survival}(r)$. The submission prints the \emph{predicted} largest safe
ratio $r^*$ \emph{before} measuring it, then reports
$|r^* - r_{\mathrm{measured}}|$.

\textbf{Measured form (N5-v2, this revision).} The degradation curve is
measured as token-level \emph{action agreement} between each compressed
checkpoint and the fine-tuned reference over a fixed prefix of the frozen
eval split (open-loop proxy; closed-loop Lyapunov rollouts are future work),
with the exact per-layer spectral deviation $\norm{W_{\mathrm{comp}} - W_{\mathrm{ref}}}_2$
(power iteration) logged alongside. The predictor uses \emph{only}
weight-space features computable at compression time, is validated
leave-one-out, and emits $r^*$ with its measured counterpart; at the current
point count the LOO $R^2$ is reported as indicative, not confirmatory.
Artifacts: \texttt{results/lyapunov/\{degradation\_seed0.json, predictor.json\}}.

\subsection{Layer 3: Statistical tail bounds on formal robustness}

Safety requirements are Signal-Temporal-Logic (STL) formulas with quantitative
robustness $\rho$ (Fainekos--Pappas semantics; $\rho>0$ = satisfied with
margin)~\cite{stl}. The challenge metric ``\% episodes in a safe equilibrium''
is the special case $\Pr[\rho \ge 0]$; we report the full $\rho$ distribution
and its tails. The scored metric is estimated empirically by the classical arms;
the Q-arm's certified bound (Sec.~\ref{sec:safety}) is its conservative
companion, stated with its direction.

% ---------------------------------------------------------------------------
\section{The Commuting-Pauli Interaction Compiler}
\label{sec:compiler}

Each cross-modal pairwise interaction term is a Pauli string
$P_a = \sigma_{a_1}\otimes\dots\otimes\sigma_{a_k}$ with coefficient $c_a$
(learned or inherited); the interaction tensor is $T = \sum_a c_a P_a$. Grouping into commuting families
and per-family Clifford diagonalization are standard VQE measurement-reduction
machinery \cite{reggio,knoclid}; our contribution is their use as an exact,
prunable compiler for cross-modal interaction tensors.
The compile pass:

\begin{table}[!t]
\caption{Compiler passes. Pass 4 is an identity, not an approximation.}
\label{tab:compile}
\centering
\small
\begin{tabular}{@{}cll@{}}
\toprule
\textbf{Pass} & \textbf{Operation} & \textbf{Output} \\
\midrule
1 & Extract interaction table $\{(c_a, P_a)\}$ & raw table \\
2 & Group into mutually commuting families $\{F_1,\dots,F_m\}$ (greedy coloring, polynomial; $m$ measured at N7) & family partition \\
3 & Per family: single Clifford $C_j$ with $C_j^\dagger P C_j \in \{I,Z\}^{\otimes k}$ & Clifford table \\
4 & Rewrite exactly $T = \sum_j C_j^\dagger D_j C_j$, $D_j$ diagonal & executable kernels \\
5 & Emit hardware table: qubits $=\lceil\log_2(\text{register})\rceil$, depth $=m$, shots $\le f(\norm{c}_1,\varepsilon)$ & Sec.~4.2 pathway \\
6 & Emit pruning certificates: dropping $F_j$ costs exactly $\norm{\sum_{a\in F_j} c_a}_1$ & ablation knob \\
\bottomrule
\end{tabular}
\end{table}

Exactness (E-comp-1): $\norm{T_{\mathrm{exact}} - T_{\mathrm{compiled}}}
\le 10^{-6}$ relative (numerics/coefficient-estimation only). Runtime cost per
family scales with the diagonal evaluation, not pairwise $O(N^2)$ enumeration
--- the attention bottleneck answered with an identity plus an $O(N\cdot k)$-style
diagonal kernel. Pruning is compositional with Eq.~\eqref{eq:safeset}: the
certified safe set survives a \emph{documented} pruning budget.

\begin{figure}[!t]
\centering
\fbox{\parbox{0.92\linewidth}{\small \centering
[Measured in N7: predicted $\norm{\sum_{a\in F_j} c_a}_1$ vs measured $\Delta$] \\[1mm]
Target: correlation $\ge 0.95$ (E-comp-2); per-family leave-one-out}}
\caption{Predicted vs measured pruning error --- the ablation column no other
submission can print.}
\label{fig:prune}
\end{figure}

\textbf{Scope and complexity, stated honestly.} The Pauli table is not
extracted from an arbitrary dense matrix (that would require up to $4^k$
terms); it is extracted from the \emph{compressed interaction accumulators} of
Sec.~\ref{sec:engine}, whose bounded rank bounds the term count, and families
beyond a top-$k$ by $\norm{c}_1$ are pruned with a certificate-bound on the
tail (R5). The grouping is greedy graph coloring --- polynomial, not optimal
--- and the realized family count $m$ is a \emph{measured} output of N7,
reported with its pruning certificate. We claim \emph{no classical
computational speedup} from the Clifford step: on classical hardware the
rewrite is an exact algebraic identity whose value is exactness, certifiable
pruning, and the hardware pathway; on quantum hardware, simultaneous
measurement in a common eigenbasis is the standard VQE shot-reduction device
that motivates commuting grouping, and k-NoCliD's non-commuting
relaxations~\cite{knoclid} mark the adjacent frontier. The quantum content is
the object language --- Pauli observables, Clifford frame, measurement basis
--- not a speed claim.

\textbf{Why this survives the QASA objection.} QASA showed a capacity-matched
classical low-rank bottleneck matches a PQC's error metrics~\cite{qasa}. Our
claim is not ``the encoding is powerful'' but ``the compiler is exact,
certifiable, and executable'': generic feature crosses have no commuting-group
structure, and Performer/FAVOR+~\cite{performer} approximates the softmax kernel
with distributional error and no notion of a family --- it cannot prune with a
printed certificate, and it emits no hardware table.

% ---------------------------------------------------------------------------
\section{Safety Evaluation Engine}
\label{sec:safety}

\subsection{Oracle}
The Sec.~5.2 seeded-scenario requirement makes the simulator a reversible
deterministic function of the seed:
\begin{equation}
O_{\mathrm{true}}: \mathrm{seed} \mapsto \mathrm{scenario} \mapsto
\mathrm{rollout}(\pi) \mapsto \rho \mapsto [\rho < 0],
\end{equation}
the object the classical arms evaluate. Amplitude estimation's end-to-end value
turns entirely on the oracle it can realize, so we state that explicitly. A
coherent oracle that evaluates a \emph{neural} rollout inside the circuit would
carry policy and dynamics as reversible subcircuits --- prohibitive depth on
hardware and exponential state-vector cost on simulators. A coherent oracle
that merely \emph{loads} $N$ pre-computed outcomes faces the $\Omega(N)$
state-preparation (QRAM) bottleneck. Our design avoids both horns:

\begin{enumerate}
\item \emph{Structured domain, free preparation.} The seeded scenario set is a
product domain (perturbation ranges per STL predicate, per the Sec.~5.2
sampling distribution), so the uniform superposition over $\log_2|S|$ seed
qubits is $H^{\otimes\log_2|S|}$ --- no amplitude-encoding of a lookup table.
Reversible arithmetic for the predicates adds $O(\log|S|)$ ancillas, so the
full AE circuit uses $O(\log|S|)$ qubits and simulated state-vector cost is
polynomial in $|S|$ ($|S|^{O(1)}$) --- tractable at our scenario counts under
the challenge's simulation mandate.
\item \emph{Closed-form certificate-margin predicates.} The Q-arm oracle
computes the \emph{certified} verdict $\tilde\rho(s) =
[\mathrm{margin}_{\mathrm{cert}}(s) < 0]$ (fusion joint 3): Layer-1 margins
are arithmetic over core norms and box geometry; Layer-2a margins are the
exported per-box SOS constants. Both are small reversible
arithmetic/comparison circuits of printed size --- never a neural rollout; the
policy is deliberately outside the circuit.
\item \emph{Conservative direction.} Sound certificates give
$\Pr[\mathrm{true\ failure}] \le \Pr[\tilde\rho(s)<0]$, so the Q-arm
estimates an upper bound on the scored quantity --- the correct ISO 26262
object, and the direction is stated on every reported number.
\end{enumerate}

Oracle circuit size (qubits, depth, gates) and the query budget are printed per
benchmark, so the claim is auditable: a \emph{query-complexity advantage under
an explicitly priced, realized oracle}. If the certificate stack degrades
(R2/R3), the cheap-predicate premise degrades with it and the classical arms
carry the section. Layer 3 is \emph{offline evaluation-time}: the $\le 100$\,ms
real-time budget governs the runtime path of Sec.~\ref{sec:arch} (compiled
diagonal kernels + monitor), which never executes an amplitude-estimation loop.

\subsection{The three-arm race}
\begin{table}[!t]
\caption{Three-arm safety estimation (pre-registered structure; complexity at
empirical $p=10^{-5}$; the Q-arm's estimand is $p_{\mathrm{cert}}$).}
\label{tab:threearm}
\centering
\small
\begin{tabular}{@{}lll@{}}
\toprule
\textbf{Arm} & \textbf{Method} & \textbf{Query cost} \\
\midrule
Q & IQAE, Bayesian posterior~\cite{iqae,tabarraei} & $\sim 10^2$--$10^3$ oracle queries $\times$ depth \\
C-strong & RESTART splitting + GEV tail fit~\cite{restart} & $10^3$--$10^4$ rollouts (measured) \\
C-naive & Monte-Carlo counting & $\sim 10^6{+}$ (cost shown, not run) \\
\bottomrule
\end{tabular}
\end{table}

Amplitude estimation estimates $p$ to absolute error $\varepsilon$ with
$O(1/\varepsilon)$ \emph{idealized-oracle} queries versus Monte-Carlo's
$O(1/\varepsilon^2)$ --- quadratic query-complexity savings on the rarest
scored quantity, standard in the QAE literature~\cite{iqae,tabarraei,mppi}. We
claim no end-to-end speedup a priori: the race measures it. The Grover depth
$1/\sqrt{p}$ is real simulator cost at our budgets --- precisely why the race
is pre-registered (R4) and the claim never assumes quantum hardware. The arms
estimate complementary quantities: Q estimates the certified (conservative)
violation probability $p_{\mathrm{cert}}$ on the cheap oracle above; C-strong
estimates the empirical failure probability from true rollouts. R4 is defined
on a common estimand: RESTART+GEV is also run on the \emph{certified}
predicate itself, so ``Q loses'' means IQAE's credible interval on
$p_{\mathrm{cert}}$ is wider than RESTART's on the same quantity at matched
total budget (printed oracle cost included); the empirical tail from true
rollouts is reported alongside as the companion number. The crossover point
(oracle complexity at which the classical arm wins) is a reported output of
N6. If Q loses, we print that; Layers 1--2 stand independently. Adjacent QAE applications include structural
reliability~\cite{tabarraei}, flow-based failure sampling~\cite{quantfp}, and
MPPI control~\cite{mppi}; Tabarraei and Das--Tanaka use lookup-table oracles
over finite ensembles --- precisely the case the oracle dilemma flags --- and
our in-domain delta is the realized cheap oracle over a structured domain fused
with exact structural certificates.

\textbf{Measured race (N6, this revision).} On the seeded heavy-tailed margin
oracle with analytic ground truth $p_{\mathrm{true}} = 1.37\times10^{-4}$
(rare-event regime, per-seed deterministic rollouts), the three arms ran at
matched query budgets: simulated Bayesian IQAE reaches a 95\% interval width
of $4.0\times10^{-5}$ with $\sim$630 oracle queries, while naive Monte-Carlo
counting needs $10^{6}$ queries for width $4.4\times10^{-5}$ (a
$\sim$1587$\times$ budget ratio at matched width); the GEV extreme-tail arm
covers the truth at $10^{6}$ queries. All covering arms are reported with
their coverage column; an arm that is narrow but misses $p_{\mathrm{true}}$
is scored \emph{worse} than a wide honest one. Two scoping statements: (i) the
RESTART \emph{splitting} half of C-strong requires rollout trajectory access,
which a seed$\to\rho$ oracle does not expose, so C-strong is implemented as
the block-minima GEV estimator with block-bootstrap CI; (ii) the Grover
amplification is \emph{simulated} --- the true amplitude lives only in the
simulator and the estimator receives measurement counts exactly as it would
on hardware. Scenario-optimization closure (Campi--Garatti, $d=2$, zero
violations): $\varepsilon \le 6.6\times10^{-4}$ at $N=10^4$
($\beta=0.01$), $\varepsilon \le 6.6\times10^{-6}$ at $N=10^6$. Conformal
coverage on fresh seeds: targets 0.99/0.95/0.90 measured as
0.9894/0.9475/0.8943. Artifacts: \texttt{results/safety/}.

\subsection{Closure}
\textbf{Scenario optimization} (Campi--Garatti): with $N$ scenarios and
solution complexity $d$, $\Pr[\text{violation}>\varepsilon]\le\beta$ whenever
$N \ge N(\varepsilon,\beta,d)$ --- printed per claim as an $(N,\varepsilon,\beta)$
triple~\cite{campi}. \textbf{Conformal calibration}: all intervals receive
split-conformal coverage correction on a held-out slice (never the training
slice; leakage rule R12). \textbf{Falsification loop}: failures found by
RESTART/IQAE feed back as training-slice augmentations, new long-tail slice
entries, and monitor training positives; the report's figure shows tail mass
shrinking across rounds.

% ---------------------------------------------------------------------------\section{Runtime Monitor}

The runtime monitor enforces the certified-safe-set at inference time: at every
step it checks that the model's predicted action token lies within the
certified ball of radius $L{\cdot}d$ around the reference action, where $L$ is
the per-layer Lipschitz product and $d$ is the input-perturbation diameter of
the scenario. If the predicted action falls outside the ball, the monitor
returns a \texttt{REFUSE\_TO\_SERVE} sentinel and logs the event; otherwise it
passes the action through. The interval logic is simple by design — its value
is that it is (i) actually enforced, not merely proposed, and (ii) formally
validated: \texttt{qicert.certify.guard.action\_in\_certified\_set} is
machine-checked against Z3 (B14g) for soundness (no out-of-ball action accepted),
completeness (no in-ball action rejected), and range enforcement, and the
guard's manifest-based load-time path verifies the self-hash of
every artifact's manifest, the hash of every listed certificate file, and
refuses service if any certified bound is NaN/Inf (B14h).

\subsection{Certified-safe-set enforcement}

At deployment, the monitor loads the signed manifest generated by
\texttt{qicert.certify.manifest} at run completion, verifies the manifest's
self-hash and each artifact's SHA-256, and then, per inference step, computes
the predicted action token and checks
\[\text{accept} \iff 0\le a < N_{\text{actions}} \;\wedge\; |a - a_{\text{ref}}| \le L{\cdot}d \;,\]where $a_{\text{ref}}$ is the certified reference action for that step (the
frozen-eval action from the same frozen episode used to derive the certificate).
The $L{\cdot}d$ radius is the product of the per-layer certificate bounds from
the QTT-compressed model (Sec.~\ref{sec:cert}, Layer~1). If the bound is
uncomputable (NaN/Inf) for any layer, the monitor refuses service --- it never
falls back to an uncomputable certificate.

The guard is deliberately minimal at the decision level; the trust comes from the
verification of the artifacts that feed it and from the formal SMT check of the
guard logic itself. A deployed monitor would additionally log each step's
$a$, $a_{\text{ref}}$, $L{\cdot}d$, and accept/reject decision to a tamper-evident
log (the manifest's self-signing pattern extends to the log).

\textbf{End-to-end demonstration (this revision).} \texttt{scripts/demo\_guard.py}
builds a certified run directory through the real certificate chain (TT-SVD
$\to$ residual compensation $\to$ sound Lipschitz bound $\to$ signed manifest)
and exercises both gates: a clean deployment is \textsc{serve}d; an
out-of-ball action is \textsc{refuse}d (safe fallback engaged); a certificate
tampered \emph{after} manifest signing is \textsc{refuse}d (hash mismatch);
a NaN bound is \textsc{refuse}d; a missing manifest is \textsc{refuse}d.
Transcript (1 serve / 4 refuse): \texttt{results/guard-demo/transcript.json}.


\label{sec:monitor}

A two-layer guard on the action stream with theorem-budgeted guarantees.

\textbf{Layer 1 --- PDU hygiene (zero extra FLOPs).} The compiled commuting-family
table (Sec.~\ref{sec:compiler}, pass 3/6) already computes a syndrome per
inference step. If the syndrome changes between consecutive steps, the packet is
gray-listed: this catches bit-flip corruption of the Pauli register --- sensor
dropouts, camera blinding, CAN-bus corruption --- the exact failure modes of the
challenge's synthetic perturbation suite.

\textbf{Layer 2 --- Classical-shadow statistics.} With $M$ random projection
directions $v_1\dots v_M$ (Huang--Kueng--Preskill shadows~\cite{shadows}),
median-of-means concentration over the scalars $s_i(x)=\langle G_i, x\rangle$
gives a statistically correct false-alarm budget at memory $O(M\cdot s)$ and a
few matvecs of latency --- continuous OOD shift that a parity bit cannot see.
A split-conformal gate sets the user-visible false-alarm rate (e.g. $\le$1\%).
Gray-list fallback: the largest-radius action from the Layer-1 certified safe
set --- the monitor never aborts into the void.

% ---------------------------------------------------------------------------
\section{Training Engine}
\label{sec:engine}

\textbf{QTT reshape} \cite{qtt}: $m=\prod m_i$, $n=\prod n_j$ (binary or
small-factor modes, bit-ordering per layer type; N2$'$ resolves the ordering).
\textbf{TT-cross} \cite{tt,ttopt}: cores $G_1\dots G_d$ are built from
$O(\sum m_i n_i r_i^2)$ maxvol queries of $W$'s entries, never forming $W$ ---
SVD-free and VRAM-safe for 7B-scale layers.

\textbf{Expressiveness, measured not assumed.} Long-range multimodal attention
is known to be low-rank-friendly (softmax kernels compress well at moderate
rank), but we do not assume it: N2 registers a per-layer rank-vs-accuracy curve
at the target ratios with a pre-registered bound (R1: $>5\%$ drop at $2\times$
demotes the compression claim), and the bond allocator then spends rank where
the \emph{certified long-tail margin} is cheapest in certificate terms.

\textbf{ALS/DMRG-native fine-tuning:} each core is optimized directly
(regularized least squares with exact multilinear gradients; AMEn enrichment
modulates bond dimension where the allocator's safety signal flags it). Memory
scales with core size, not layer size. Training passes: P1 baseline fine-tune;
P2 compress $\rightarrow$ ALS fine-tune (the deliverable); P3 safety-critical
last-pass on the RESTART/IQAE-falsified failure slice only (tail-shrinkage
evidence, $\sim$hundreds of hard examples).

\textbf{Safety-budgeted bond allocator:} per layer,
\begin{equation}
r^* \;=\; \arg\min_r \max_{\text{slices}} \big( \text{certified long-tail margin at } r \big),
\end{equation}
using (1) layer sensitivity to the certified-margin computation, (2) bond-spectrum
entropy, (3) monitor hygiene bits. The long-tail slice is a design selector, not
a report line.

% ---------------------------------------------------------------------------
\section{Experimental Setup}
\label{sec:setup}

\textbf{Tracks and baselines (Sec.~5.4):} AD --- LLaVA-1.5-7B (Alpamayo-adjacent
head) on nuScenes \cite{nuScenes}/CARLA-style slices, vs. INT8 bitsandbytes;
robotics --- MiniVLA-class backbone (MIT code, HF checkpoints; Open
X-Embodiment lineage \cite{oxe}) on LIBERO-style slices, vs.
INT8. Strong extra baselines: SVD-at-matched-ratio, tuned Performer at matched
capacity, RESTART+GEV (not naive MC). \textbf{Hardware:} Kaggle 2$\times$T4
(30 GPU-h/wk, 120 h ceiling), local RTX 5060 for small runs; all budgets are
quota-hours, silicon-hours logged separately. \textbf{Discipline:} mean $\pm$ std
over $\ge$3 seeds; every report table has a ledger row
(\texttt{research/experiments.md}); every experiment has a pre-registered kill
criterion.

\begin{table}[!t]
\caption{Experiment matrix (pre-registered; GPU-h are quota-hours).}
\label{tab:exp}
\centering
\small
\begin{tabular}{@{}cllc@{}}
\toprule
\textbf{ID} & \textbf{Experiment} & \textbf{Pillar} & \textbf{GPU-h} \\
\midrule
N1 & AD backbone fine-tune + INT8 reference & A & 8 \\
N2 & $\{$TT,QTT$\}\times$TT-cross sweep, 6 bond plans $\times$ 2 backbones & A & 24 \\
N2$'$ & Bit-ordering sensitivity (3 orders, 1 layer, 1 seed) & A & 2 \\
N3 & Layer-1 certificate table (Lipschitz products + tightness ratio $\kappa$) & A & 2 \\
N4 & Layer-2a SOS boxes (3 points $\times$ 2 tracks $\times$ slices) & A & 3 \\
N5 & Layer-2b Lyapunov curve + bond-spectrum predictor & A & 6 \\
N6 & Layer-3 safety suite (IQAE vs RESTART+GEV vs MC; oracle cost + crossover) & C & 8 \\
N7 & Compiler exactness + prune certificates + latency & B & 6 \\
N8 & (optional) PQC RL toy head + noise sweep & tertiary & 4 \\
N10 & Syndrome--shadow monitor demo & monitor & 3 \\
N11 & ALS/DMRG-native fine-tune, both backbones & engine & 8 \\
N12 & Safety-critical last-pass on falsified slice & engine & 3 \\
N13 & Cross-track mirror (N2/N3/N5 on second backbone) & A & 12 \\
N14 & Clean-env end-to-end bench audit & all & 4 \\
\bottomrule
\end{tabular}
\end{table}

\textbf{Budget rollup:} 87 GPU-h core, 33 h headroom vs. the 120 h ceiling;
weeks 1--6 schedule in \texttt{08-experiments.md}. Week 1: N1+N2$'$+N3
(the cheapest exposure of the two load-bearing facts: does the baseline behave,
and does QTT compress with certificates attached?). The optional N8 PQC demo
uses a quantum natural policy gradient head \cite{qnpg} for illustration only;
it is cut first under quota pressure (R10).

% ---------------------------------------------------------------------------
\section{Results Framework}
\label{sec:results}

Per challenge Sec.~5.5, all values are mean $\pm$ std over $\ge$3 runs. Tables
below state the pre-registered targets; measured values land from the ledger
and are inserted verbatim (no rounding-into-a-target). No number in this report
is fabricated: each row is regenerable from the per-run \texttt{run.json} in
\texttt{results/} plus the ledger (\texttt{results/ledger.csv}).

\textbf{Measured so far (this revision):} the classical-energy N1v2 baseline
(MiniVLA-1B, frozen 108-episode LIBERO spatial split, streamed eval, seed~0 of
3 seeds, 1200 LoRA-finetune steps at batch 2, local RTX 5060):
\begin{center}
\small
\begin{tabular}{@{}lc@{}}
\toprule
\textbf{Metric} & \textbf{Seed 0 (mean $\pm$ 95\% Wilson CI over 6{,}496 batches)} \\
\midrule
Fine-tuned action accuracy & 0.4453 $\pm$ 0.0029 $[$0.4424, 0.4482$]$ \\
INT8-of-fine-tuned (torchao Int8WeightOnlyConfig, GPU) & 0.0000 $[$0.0000, 0.0000$]$ \\
INT8 vs FT drop (matched protocol, both legs, full split) & $-$0.4453 \\
\bottomrule
\end{tabular}
\end{center}
A zero INT8-of-fine-tuned value on this fresh LoRA-fine-tuned model is expected
and honest: naive weight-only INT8 without calibration collapses a
fp16-fine-tuned qwen2.5-0.5B's action-token logits. The honest comparison for
N2 is therefore \emph{N2$'$-compressed (INT4 base + certified TT residual) vs.
calibrated INT8 at matched compression ratio}, not naive weight-only INT8 --- the
N2$'$ pipeline includes a Hessian-aware residual-fitting stage (GPTQ-intrinsic
recipe + the QuaSAR-stable-pseudoinverse guard) that is the actual
``INT8-competitive'' arm. Seeds 1 and 2 of the N1v2 baseline were launched
concurrently with this writing and will land in the ledger as they finish.

\textbf{Measured compression status (full-split, validated harness, this
revision).} The eval harness was independently validated: loading the
unmodified fine-tuned weights through the compressed-eval swap path
reproduces the streamed reference to $+0.0000$ (0.4468, 6{,}496 batches,
114{,}497 action tokens). On that protocol, per-channel \emph{calibrated}
INT8 is near-lossless at 2.00$\times$ (0.4466, $\Delta=-0.0001$) --- which
also proves the earlier 0.0000 naive-INT8 row was a calibration artifact,
not a property of INT8 --- while the certified TT+residual track at its
best measured point holds 0.1529 at 2.86$\times$ ($\Delta=-0.2939$), with
all 168 per-layer certificates sound. The pre-registered R1 bar (some point
at $\ge 2\times$ with $\Delta \ge -0.05$) is therefore \textbf{not met} by
the uniform and Frobenius-residual arms; the activation-weighted fit
(N2R-v2, GPTQ-style) and mixed allocation are the pre-registered next
measurements, and the honest current claim is the certificate stack (the
only compressed model with sound runtime certificates), not accuracy
leadership. Every number in this paragraph regenerates from
\texttt{results/} run artifacts + the ledger.

\begin{table}[!t]
\caption{Main results table (structure; $\dagger$ = pre-registered target).}
\label{tab:main}
\centering
\small
\begin{tabular}{@{}lccc@{}}
\toprule
\textbf{Row} & \textbf{qicert-QTT} & \textbf{INT8} & \textbf{SVD} \\
\midrule
Accuracy @ 2$\times$ (drop vs dense) & $\le$5\%$^{\dagger}$ & measured & measured \\
Accuracy @ 4$\times$ (drop vs dense) & target $\le$5\%$^{\dagger}$ & measured & measured \\
Certified-safe-set \% @ 2$\times$ & $>$90\%$^{\dagger}$ & --- (not computable) & --- \\
Certified-safe-set \% @ 4$\times$ & target $>$50\%$^{\dagger}$ & --- & --- \\
Latency (ms, stated profile) & $\le$100 & measured & measured \\
\bottomrule
\end{tabular}
\end{table}

\begin{table}[!t]
\caption{Compiler acceptance (E-comp) and three-arm safety (per slice).}
\label{tab:compacc}
\centering
\small
\begin{tabular}{@{}ll@{}}
\toprule
\textbf{E-comp} & \textbf{Target} \\
\midrule
1 exactness & $\norm{T_{ex}-T_{comp}}\le 10^{-6}$ rel. \\
2 prune-certificate corr. & predicted vs measured $\ge 0.95$ \\
3 latency & $\le$100 ms at full table \\
4 hardware self-consistency & compiled vs simulated within tolerance \\
\bottomrule
\end{tabular}
\vspace{2mm}
\begin{tabular}{@{}ll@{}}
\toprule
\textbf{Arm} & \textbf{Output} \\
\midrule
Q (IQAE) & credible interval on $p_{\mathrm{cert}}=\Pr[\mathrm{margin}_{\mathrm{cert}}<0]$ (conservative) \\
C-strong (RESTART+GEV) & profile-likelihood CI on $p$ \\
C-naive (MC) & cost shown; Clopper--Pearson CI \\
Scenario-opt & $(N,\varepsilon,\beta)$ triple \\
Conformal & coverage on held-out slice \\
\bottomrule
\end{tabular}
\end{table}

% ---------------------------------------------------------------------------
\section{Ablation and Quantum Justification}
\label{sec:ablation}

Per challenge Sec.~5.5, the QI-isolating ablation is mandatory and is designed
to be the strongest section of the report.

\textbf{A1 (compression$\times$certificate):} at matched ratio, INT8 matches
accuracy but cannot report the certificate column --- the ablation shows the
asymmetry directly (the number that cannot be computed is the win). The
tightness ratio $\kappa$ is reported per point: the certificate column is
conservative-but-quantified, and Layer~2a closes the gap locally.

\textbf{A2 (compiler vs encoding):} per-family leave-one-out with \emph{predicted}
error $\norm{\Sigma_{a\in F_j} c_a}_1$ vs measured $\Delta$; and compiler-vs-
encoding: the exact identity vs a bare feature cross (shows the family structure
is what buys exactness, not the encoding label). Tuned Performer at matched
capacity is the classical ceiling; winning it is upside, losing it does not kill
the claim (the primary axis is exactness + certificates + hardware table).

\textbf{A3 (IQAE vs strong classical):} the three-arm race at matched query
budget; if Q loses, the report prints that (R4) and the claim rests on Layers
1--2 and the strong classical arm.

\textbf{A4 (monitor):} syndrome--shadow vs autoencoder baseline: correction rate
at matched overhead; false-alarm rate vs the shadow theorem bound ($\le$3$\times$
or demoted, R7).

\textbf{A5 (training engine):} ALS-in-compressed vs dense-then-compress at
matched ratio on the safety slice ($\le$2\% worse or demoted, R6).

% ---------------------------------------------------------------------------
\section{Limitations and Risks}
\label{sec:limitations}

Pre-registered kill/demote criteria (full register in
\texttt{09-risks-kills.md}): R1 compression accuracy $>$5\% drop at 2$\times$
$\rightarrow$ CompactifAI-style softer claim + certificate story; R2 Layer-1
bound too loose ($<$50\% certified safe set) $\rightarrow$ SOS becomes primary;
R3 SOS SDP intractable $\rightarrow$ Lyapunov curve primary, SOS reported as
``attempted, too slow''; R4 IQAE loses to RESTART $\rightarrow$ RESTART stands,
IQAE = future work; R5 family count too large $\rightarrow$ top-$k$ by
$\norm{c}_1$ with certificate-bound; R6 ALS-in-compressed $>$2\% worse
$\rightarrow$ demote to post-processing; R7 shadow FPR $>$3$\times$ theorem
$\rightarrow$ PDU hygiene only; R8 allocator tradeoff $\rightarrow$ reported
honestly; R9 cross-track non-comparable $\rightarrow$ pipeline claim only;
R10 quota overrun $\rightarrow$ N8 cut first, then N13 reduced; R11 clean-env
failure $\rightarrow$ report held; R12 conformal leakage $\rightarrow$ claim
degraded until re-run.

\textbf{Explicit framing disclaimers.} We claim no classical speedup from the
Clifford step (Sec.~\ref{sec:compiler}) and place amplitude estimation strictly
offline with printed oracle costs (Sec.~\ref{sec:safety}); both are stated
in-body so no stronger claim can be imputed. No end-to-end advantage over
classical estimators is claimed for IQAE a priori --- the pre-registered race
at matched total budget and its measured crossover decide (Sec.~\ref{sec:safety}).
Runtime parity with classical baselines is an acceptance target, not an
assumption: E-comp-3 ($\le 100$\,ms at the full compiled table) and the N14
clean-env bench are the measured evidence.

\textbf{Three things that cannot kill the core claim:} (1) the exact
Lipschitz-product certificate --- a theorem, not a benchmark; (2)
compilation-as-identity + pruning certificate --- a mechanism, not a number;
(3) the three-layer safety methodology --- a formal spec plus raced estimators
plus conformal closure, not a score. The final table is built from whatever
layers survive their criteria; the package always runs and reports something
certified.

\textbf{What is actually measured vs. pre-registered in this revision:} the
classical-energy N1v2 baseline (MiniVLA-1B, full frozen 108-episode LIBERO
spatial split, stream-eval, 3 seeds $\times$ 1200 LoRA fine-tune steps at
batch 2 on local RTX 5060; seed 0 is complete with fine-tuned accuracy
0.4468 (Wilson 95\%, streamed over 6{,}496
batches) --- see \texttt{results/N1v2}; the INT8-of-fine-tuned leg runs on
\texttt{torchao Int8WeightOnlyConfig} (GPU) over the same 6{,}496 frozen
batches, giving a matched-protocol $\Delta$(INT8 $-$ FT) per seed once seeds
1--2 complete (staged one at a time; \texttt{EXECUTION-PLAN.md}). The
measured table now also includes the FT-anchor validation ($+0.0000$), the
calibrated-INT8 comparator (0.4466 at 2.00$\times$), and the certified
TT+residual verdicts (0.1529 at 2.86$\times$, 168/168 certificates sound) ---
see the Results Framework measured-status paragraph. No number in this revision is fabricated; each final table row
is regenerable from \texttt{results/N1v2/\{run\_id\}/run.json} + the ledger.

\textbf{Certification framing:} the STL specs map to ISO 26262 / IEC 62061 /
ISO 10218 artifacts (hazard $\rightarrow$ safety goal $\rightarrow$ STL formula
$\rightarrow$ robustness predicate $\rightarrow$ certified tail probability),
and the certificate stack provides the evidence trail an assessor expects.

\textbf{Scope generalization (explicit limitation).} All compression
measurements to date are on one backbone (MiniVLA-1B, qwen2.5-0.5B LLM) and
one task suite (LIBERO spatial, frozen 108-episode split). The measured
TT-truncation collapse at $\ge 2\times$ is partly a property of this
backbone's unusually flat weight spectra (information spread across singular
directions; little discardable redundancy) and its fragile action head;
literature reports structured and transformer backbones compressing far more
gracefully. The certificate machinery, the guard, the estimator race, and the
conformal closure are \emph{model-agnostic by construction} (they operate on
whatever weights the pipeline produces). Phase-2 work: repeat the full
pipeline on a second backbone and a second task suite; extend Layer-2b to
closed-loop Lyapunov rollouts; add an INT4 comparator for the $>2\times$
regime; and a hardware-queue run. Layer-2a SOS boxes remain scoped behind
the Julia pin (Q22) and are reported as attempted-not-run at submission if
the pin does not land.

% ---------------------------------------------------------------------------
\section{Resource Declaration}

\textbf{Hardware (quantum simulation tier):} Kaggle 2$\times$NVIDIA T4 16\,GB (primary), local RTX 5060 laptop (secondary, development + classical baseline).  
\textbf{Hardware (classical-energy baseline tier):} local RTX 5060 (1$\,\times\,$8\,GB) — the N1/N2 classical-energy runs that validate the energy accounting and provide the QNN-vs-classifier fidelity context for Q3, 
all on this host only (no GPU cluster relied upon).  
\textbf{Total:} 87 GPU-h core (quota-hours) + headroom; silicon hours logged separately in the ledger.  
\textbf{Simulation stack (quantum):} CUDA-Q (pinned version), PennyLane, tntorch --- versions in \texttt{environment.yml}; CUDA-Q available through Docker on this host.  
\textbf{Energy baseline stack (classical):} PyTorch 2.14 + torchao + peft + HF transformers 5.x + Min\'iVLA-1B backbone weights --- versions in \texttt{docs/resource-declaration.md} (auto-generated from run provenance during the live bench).  
\textbf{Streamed eval workload note:} on the classical-energy baseline tier, the eval split is consumed lazily (two fresh stream iterators, one per eval leg) so the whole 108-episode held-out set (6{,}496 batches, ~16\,GB of batches in the earlier design) fits stably in ~10\,MB per batch without materialization --- the earlier materialized-batch design OOM-killed processes silently; the streaming fix is in \texttt{bench/compress.py} since 2026-09-11.
\textbf{Shot/gate budgets:} per quantum benchmark, stated in the compiler's pass-5 table.  
\textbf{Energy:} measured by the \texttt{qicert} profiler (kWh + CO$_2$e at declared grid intensity). Reproducibility evidence: \texttt{qicert.bench.all} from a fresh virtualenv, verified by the certify suite (B14a--B14h) and --skip-int8 handoff for N2.

% ---------------------------------------------------------------------------
\section{Trust Augmentation: Making Every Number Defensible}
\label{sec:trust}

The single biggest failure mode in this kind of submission is an implicit or
explicit claim that is not traceable to a defensible number --- a number that
a skeptical reader could not reproduce without asking a follow-up question. We
mitigate this with a \emph{trust-augmentation stack} (B14 family) that makes
every reported number one of three things: measured with a confidence interval,
checked against an independent verifier, or witnessed by a machine-checkable
artifact. \textbf{No number in the final tables is claimed on the authority of
the person who ran it.} The stack is pure-Python (no external formal-methods
heavy-lifting required) and is developed and tested in this repository:

\begin{itemize}
\item \textbf{B14a --- certifying-algorithm witnesses + property-based tests
(Hypothesis).} Every certificate artifact ships as a \texttt{Witness} dataclass
containing the serialized cores, the claimed Lipschitz float, and the
bound-kind. An \emph{independent} verifier re-derives the bound from the cores
and asserts equality to tolerance. Property tests assert gauge-invariance of the
TT contraction (bond rescaling with explicit inverse on a chain leaves the
contracted operator invariant to machine precision), product-of-core-norms
determinism, power-iteration monotonicity, \emph{and} tamper detection: a
witness whose shipped bound has been perturbed away from the true value must
\texttt{FAIL} the verifier --- the witness is an integrity artifact, not a
decoration. (See \texttt{tests/test\_certify\_properties.py}.)

\item \textbf{B14b --- statistical robustness (Clopper--Pearson).} The safety
layer's failure probability is a tail event ($p_{\text{cert}} = \Pr[\text{margin}<0]$)
and must be reported with a confidence interval, not a point estimate. We use
two-sided Clopper--Pearson binomial intervals (the honest, conservative interval
regulators expect) over the frozen 108-episode held-out split. Property tests
verify interval containment, boundary behavior, and monotonicity in the
violation count --- so the reported interval is guaranteed to be at least as
conservative as claimed. (See \texttt{tests/test\_certify\_properties.py}.)

\item \textbf{B14c --- interval/exact-arithmetic certificate cross-check.} The
float pipeline's certificate bound for a small layer is independently recomputed
in 50-digit \texttt{mpmath} precision (power-iteration on $A^TA$ --- dependency-
free; this version fixes an earlier prototype that called the non-existent
\texttt{mp.\,linalg.\,svdvals}) and, for small test layers, in exact
\texttt{fractions.\,Fraction}. The float bound is accepted only if it matches
the 50-digit recomputation to tolerance and, where conversion-safe, the exact
fraction. This catches ``your bound is a float artifact'' attacks. (See
\texttt{tests/test\_certify\_properties.py} + the
\texttt{qicert.certify.cross\_check} module.)

\item \textbf{B14g --- SMT falsification of the runtime guard (Z3).} Rather than
proving the guard sound by hand, we ask Z3 to \emph{disprove} its properties: the
SMT encoding of \texttt{action\_in\_certified\_set} is checked for
soundness (no out-of-ball action accepted), completeness (no in-ball action
rejected), and range enforcement. If Z3 found a counterexample, the guard would
\emph{not} be deployable; the unsat results in
\texttt{tests/test\_certify\_guard.py} are the machine-validated assurance the
guard is sound and complete on its stated domain. We also cross-validate the
implementation against an independent Z3 model on 200 random
(action,reference,margin,n) points.

\item \textbf{B14h --- signed manifests + runtime assertion guard.} At run
completion, \texttt{qicert.certify.manifest.from\_dir} writes a manifest that
SHA-256-hashes every artifact (model weights, certificates, frozen eval split,
provenance files) and self-signs by hashing the manifest itself. The runtime
guard \texttt{qicert.certify.guard.check\_certs\_before\_serve} re-verifies
the manifest self-hash on load (tamper detection), verifies every listed
artifact's hash, and refuses service if any bound is NaN/Inf. Property tests
verify tamper detection on real files (mutated artifact, forked manifest entry,
missing artifact, NaN bound) --- so the ``receipts are enforced'' narrative is
real rather than aspirational.
\end{itemize}

\textbf{The honest story.} With the B14 stack, every number in the final tables
says: \emph{this was measured with a confidence interval / checked against an
independent verifier / witnessed by a machine-checkable artifact}, and the
composition rule itself (the runtime guard) is machine-verified by a third-party
SMT solver. That is stronger than most submissions' entire verification story.
The B14 stack is in development in this repository and is being refined as the
quantum experiments land; it is architecture-complete (all of B14a--B14h
implemented and tested) even where the quantum artifacts it validates are still
pending.

% ---------------------------------------------------------------------------
\section{Conclusion}

We submit objects the baselines structurally cannot produce: exact certificates
that are functions of the compressed cores, an exact compilation identity with
per-family provable pruning, quadratic query-complexity savings on the rarest
scored quantity under a realized, priced oracle, and a package that reproduces
every table from a clean environment.
All claims are either theorems, structural properties, or pre-registered
measurements with kill criteria --- the honesty contract of Sec.~5.5
operationalized. The full argument, theory, and risk register are maintained in
the submission folder of the \texttt{qicert} repository.

% ---------------------------------------------------------------------------
\begin{thebibliography}{99}

\bibitem{alpamayo} NVIDIA and Y.~Wang et al., ``Alpamayo-R1: Bridging Reasoning
and Action Prediction for Generalizable Autonomous Driving in the Long Tail,''
arXiv:2511.00088, 2025.

\bibitem{pi0} K.~Black et al., ``$\pi_0$: A Vision-Language-Action Flow Model
for General Robot Control,'' arXiv:2410.24164, 2024.

\bibitem{compactifai} A.~Tomut et al., ``CompactifAI: Extreme Compression of
Large Language Models using Quantum-Inspired Tensor Networks,''
arXiv:2401.14109, 2024.

\bibitem{liu} J.~Liu et al., ``Towards Provably Efficient Quantum Algorithms
for Large-Scale Machine-Learning Models,'' Nature Communications 15(1):434,
2024.

\bibitem{tt} I.~V. Oseledets, ``Tensor-Train Decomposition,'' SIAM J. Sci.
Comput. 33(5):2295--2317, 2011.

\bibitem{qtt} I.~V. Oseledets, ``Approximation of $2^d\times2^d$ matrices using
tensor products,'' arXiv:1007.1170, 2010.

\bibitem{ttopt} M.~Boussé et al., ``TTOpt: A MATLAB package for Tensor-Train
optimization,'' arXiv:2205.00293, 2022.

\bibitem{sos-tt} L.~Balada Gaggioli, D.~Henrion, and M.~Korda, ``Composition
and tensor train structure in polynomial optimization,'' arXiv:2604.17563,
2026.

\bibitem{qasa} S.~Chen and Y.~Kuo, ``Quantum-inspired attention architectures,''
Quantum Science and Technology 11:035051, 2026.

\bibitem{performer} K.~Choromanski et al., ``Rethinking Attention with
Performers,'' arXiv:2009.14794, 2020.

\bibitem{reggio} A.~Reggio et al., ``Pauli grouping for measurement
reduction,'' arXiv:2305.11847, 2023.

\bibitem{knoclid} N.~P.~D. Sawaya, D.~Camps, B.~DalFavero, N.~M.~Tubman,
G.~M.~Rotskoff, and R.~LaRose, ``Non-Clifford diagonalization for measurement
shot reduction in quantum expectation value estimation,'' arXiv:2408.11898,
2024 (v3 2025).

\bibitem{stl} G.~E. Fainekos and G.~J. Pappas, ``Robustness of temporal logic
specifications,'' in Proc. Formal Methods and Models for Co-Design (MEMOCODE),
2009.

\bibitem{iqae} D.~Grinko, J.~Gacon, C.~Zoufal, and S.~Woerner, ``Iterative
Quantum Amplitude Estimation,'' npj Quantum Information 7:52, 2021;
arXiv:1912.05546.

\bibitem{tabarraei} A.~Tabarraei, ``Bayesian sequential quantum amplitude
estimation for rare-event structural failure probability,'' arXiv:2607.18996,
2026.

\bibitem{quantfp} A.~I. Weinberg, ``QuantFPFlow: Quantum amplitude estimation
for Fokker--Planck policy optimisation in continuous reinforcement learning,''
arXiv:2605.16429, 2026.

\bibitem{mppi} G.~Das and T.~Tanaka, ``Model predictive path integral control
as a quantum query problem,'' arXiv:2607.28851, 2026.

\bibitem{restart} M.~Vill\'en-Altamirano and J.~Vill\'en-Altamirano,
``RESTART: A method for accelerating rare event simulations,'' in Proc. ITC-13,
1991.

\bibitem{campi} M.~C. Campi and S.~Garatti, ``The exact feasibility of
randomized solutions of uncertain convex programs,'' SIAM J. Optim. 19(3),
2008.

\bibitem{shadows} H.-Y. Huang, R. Kueng, and J. Preskill, ``Predicting many
properties of a quantum system from very few measurements,'' Nature Physics
16:1050--1057, 2020.

\bibitem{nuScenes} H.~Caesar et al., ``nuScenes: A Multimodal Dataset for
Autonomous Driving,'' CVPR, 2020.

\bibitem{oxe} Open X-Embodiment Collaboration, ``Open X-Embodiment: Robotic
Learning Datasets and RT-X Models,'' arXiv:2310.08864, 2023.

\bibitem{qnpg} J.~Meyer et al., ``Quantum Natural Policy Gradients,''
arXiv:2304.13571, 2023.

\end{thebibliography}

\end{document}
